\section{Statement of Compliance}

This trial will be conducted in accordance with the protocol, the principles of
the International Council for Harmonisation (ICH) Good Clinical Practice (GCP)
guideline E6(R3), and all applicable United States regulations governing both
the investigational drug and the investigational device, together with the
Physical AI overlay that governs the on-premises large language model (LLM)
directed robotic system. Because this is a combined Investigational New Drug
(IND) and Investigational Device Exemption (IDE) first-in-human study, two
distinct but harmonized regulatory spines apply concurrently; both are carried
in full below, and neither is subordinated to the other. The Sponsor-Investigator
attests that the conduct of every study activity, the protection of every
enrolled participant, and the integrity of every recorded datum will satisfy the
more stringent of the two spines wherever they differ.

\subsection{Drug Arm Compliance (IND, 21 CFR Part 312)}

The perioperative daraxonrasib (RMC-6236) component is conducted under an IND in
accordance with Title 21 of the Code of Federal Regulations. Specifically, the
study complies with 21 CFR part 50 (Protection of Human Subjects, including the
informed-consent requirements of \S50.20 through \S50.27); 21 CFR part 54
(Financial Disclosure by Clinical Investigators); 21 CFR part 56 (Institutional
Review Boards); and 21 CFR part 312 (Investigational New Drug Application),
including the sponsor and investigator responsibilities of subparts B and D, the
safety-reporting requirements of \S312.32, and the IND-application and amendment
provisions of \S312.20 through \S312.31. The drug arm is structured as a Phase 1,
3+3 dose-finding evaluation; the IND is in effect following the 30-day review
period of \S312.40, and no participant is dosed before the IND is in effect.

\subsection{Device Arm Compliance (IDE, 21 CFR Part 812)}

The on-premises LLM-directed eight-arm robotic pancreaticoduodenectomy (Whipple)
system is an investigational medical device evaluated under an IDE in accordance
with 21 CFR part 812. The Sponsor-Investigator and the reviewing Institutional
Review Board (IRB) have determined that the device presents a significant risk
within the meaning of \S812.3(m), because it is a surgical-intervention device
whose patient-contact motion, force application, vascular-exclusion gating, and
emergency-stop behavior present a potential for serious risk to the health,
safety, or welfare of a participant. The study therefore proceeds only under an
IDE approved by both the FDA and the IRB. The device arm additionally complies
with 21 CFR part 50 and 21 CFR part 56 for human-subject protection and IRB
oversight, and with 21 CFR part 11 (Electronic Records; Electronic Signatures)
for the validation, audit-trail, and electronic-signature controls that
underpin the hash-chained command and telemetry record described in \S10. The
device component is framed as an early feasibility study consistent with FDA
guidance for significant-risk devices, in which limited clinical experience is
sought because nonclinical simulation alone cannot supply the information needed
to advance development.

\subsection{Physical AI Overlay (21 CFR Part 312, Subpart J)}

The autonomy-bearing behavior of the LLM-directed system is governed by the
Physical AI overlay codified in 21 CFR part 312, Subpart J (\S312.400 through
\S312.405), as adapted for end-to-end Physical AI oncology investigations
\cite{kawchak2026adapt312}. Subpart J requires, as preconditions to first-in-human
use, a Phase 0 simulation-validation gate (at least two independent simulation
frameworks reproducing motion within a $\leq 2$ mm positional and $\leq 0.5$ N
force tolerance), a documented Unified Safety Level (USL) readiness rating of
$\geq 7.0$ for surgical deployment, and classification of the system as a Class
II collaborative device under continuous human oversight with a guaranteed
emergency-stop latency of $\leq 500$ ms. These overlay requirements operate in
addition to, and never in place of, the part 312 and part 812 obligations above.
The hash-chained audit trail required under \S312.404 satisfies the part 11
electronic-records standard and renders every LLM advisory, every robot command,
and every safety-limit event re-traceable to a deterministic seed and a specific
repository commit.

\subsection{Good Clinical Practice and Electronic Records}

The trial adheres to ICH E6(R3) GCP, the international ethical and scientific
quality standard for designing, conducting, recording, and reporting clinical
investigations involving human participants. All electronic clinical data,
source records, and the cyber-physical telemetry stream are maintained under 21
CFR part 11, with validated systems, secure time-stamped audit trails that do not
obscure prior entries, role-based access controls, and electronic signatures
bound to the records they authenticate. The hash-chained record specified by the
Physical AI overlay provides cryptographic continuity across the screening,
operative, and follow-up phases.

\subsection{Trial Registration}

The study will be registered on ClinicalTrials.gov before enrollment of the
first participant. It is registered as an Interventional study. The drug
component is classified as Phase 1 (dose-finding for daraxonrasib), while the
device component carries Phase: Not Applicable, consistent with the
ClinicalTrials.gov convention that FDA-defined phase categories apply to drugs
and biologics rather than to devices. The Primary Purpose is recorded as Device
Feasibility together with the dose-finding objective of the drug arm. The study
is open-label and single-arm, with staggered (sentinel) enrollment and
prespecified pause and stopping rules.

\subsection{Statement of Compliance Pathway}

Figure 1 renders the combined IND and IDE compliance spine with the Physical AI
Subpart J overlay. The daraxonrasib drug arm (IND, part 312) and the robotic
Whipple device arm (IDE, part 812, significant-risk) converge through the
Phase 0 simulation-validation gate, the USL readiness verification, and the
Class II collaborative classification into a single first-in-human protocol.

\begin{mermaidfig}
\node[mmin]  (pre)   at (0,0)       {Pre-IND meeting\\\S312.47 plus\\Physical AI review};
\node[mmstep](ind)   at (0,-2.4)    {IND in effect (30-day)\\daraxonrasib\\Phase 1, 3+3};
\node[mmin]  (sr)    at (4.6,0)      {Significant-risk\\determination\\\S812.3(m)};
\node[mmstep](ide)   at (4.6,-2.4)   {IDE approved\\(FDA plus IRB)\\early feasibility};
\node[mmstep](p0)    at (2.3,-4.8)   {Phase 0 sim validation\\$\geq 2$ frameworks\\$<2$ mm / $<0.5$ N};
\node[mmstep](usl)   at (2.3,-7.0)   {USL readiness $\geq 7.0$\\pre-procedure\\safety matrix};
\node[mmgoal](cls)   at (2.3,-9.2)   {Class II collaborative\\continuous oversight\\$\leq 500$ ms e-stop};
\node[mmgoal](enr)   at (2.3,-11.4)  {First-in-human\\enrollment\\combined IND plus IDE};
\draw[mmedge] (pre) -- (ind);
\draw[mmedge] (sr)  -- (ide);
\draw[mmedge] (ind.south) |- (p0.west);
\draw[mmedge] (ide.south) |- (p0.east);
\draw[mmedge] (p0)  -- (usl);
\draw[mmedge] (usl) -- (cls);
\draw[mmedge] (cls) -- (enr);
\end{mermaidfig}
\figcaption{Figure 1. Combined IND / IDE regulatory pathway with the Physical AI
overlay. The daraxonrasib drug arm (IND, part 312) and the robotic Whipple device
arm (IDE, part 812, significant-risk per \S812.3(m)) converge through the
Subpart J overlay of Phase 0 simulation validation, USL readiness, and Class II
collaborative classification into one first-in-human protocol.}

\subsection{Signatures}

The Sponsor-Investigator has reviewed and approved this protocol and confirms
that the study will be conducted in compliance with the protocol, ICH E6(R3) GCP,
and the federal regulations cited above. By signature below, each signatory
agrees to conduct or supervise the study in accordance with the protocol and all
applicable requirements.

\vspace{0.6em}

\noindent\begin{tabularx}{\textwidth}{L{8.2cm} Y}
\toprule
\textbf{Sponsor-Investigator} & \textbf{Site Principal Investigator} \\
\midrule
\vspace{1.6em}\rule{6.0cm}{0.4pt} & \vspace{1.6em}\rule{6.0cm}{0.4pt} \\
Kevin Kawchak & Name (printed) \\
CEO, ChemicalQDevice & Title / Institution \\[0.8em]
Signature and date: \rule{2.6cm}{0.4pt} & Signature and date: \rule{2.6cm}{0.4pt} \\
\bottomrule
\end{tabularx}

\vspace{0.4em}

\noindent The signatory acknowledges that this is an independent research draft;
it is not an active IND or IDE and is not endorsed by the FDA, NIH, HHS, an IRB,
ICH, or any sponsor.
